% Chunk: \S5.6 General Irreducible Representations of the Homogeneous Lorentz Group
% Source: Weinberg QFT Vol. I, physical pp. 252--256 (printed pp. 229--233; begins and ends mid-page)
% Status: source-reviewed, compile-clean, and visually verified; equations (5.6.1)--(5.6.20); no figures; no uncertainties

\subsection[General Irreducible Representations of the Homogeneous Lorentz Group]{General Irreducible Representations of the Homogeneous Lorentz
Group\texorpdfstring{\protect\footnotemark}{}}
\footnotetext{This section lies somewhat out of the book's main line of
development, and may be omitted in a first reading.}

We shall now generalize from the special cases of vector and Dirac fields to
the case of a field that transforms according to a general irreducible
representation of the homogeneous Lorentz group. All fields may be constructed
as direct sums of these irreducible fields.

A general representation of the proper orthochronous homogeneous Lorentz group
(or, more properly, its infinitesimal part) is provided by a set of matrices
$\mathcal{J}_{\mu\nu}$ satisfying the same commutation relations \eqref{eq:5.4.4} as the
generators of the group
\begin{equation}
[\mathcal{J}_{\mu\nu},\mathcal{J}_{\rho\sigma}]
=i(\mathcal{J}_{\rho\nu}\eta_{\sigma\mu}
+\mathcal{J}_{\mu\sigma}\eta_{\nu\rho}
-\mathcal{J}_{\sigma\nu}\eta_{\rho\mu}
-\mathcal{J}_{\mu\rho}\eta_{\nu\sigma}).
\tag{5.6.1}\label{eq:5.6.1}
\end{equation}
(Of course, $\mathcal{J}_{\mu\nu}=-\mathcal{J}_{\nu\mu}$, and indices on
$\mathcal{J}_{\mu\nu}$ are as usual raised or lowered by contraction with
$\eta^{\mu\nu}$ or $\eta_{\mu\nu}$.) To see how to construct such matrices,
first divide the six independent components of $\mathcal{J}_{\mu\nu}$ into two
three-vectors: an angular momentum matrix
\begin{equation}
\mathcal{J}_1=\mathcal{J}_{23},\qquad
\mathcal{J}_2=\mathcal{J}_{31},\qquad
\mathcal{J}_3=\mathcal{J}_{12}
\tag{5.6.2}\label{eq:5.6.2}
\end{equation}
and a boost
\begin{equation}
\mathcal{K}_1=\mathcal{J}_{10},\qquad
\mathcal{K}_2=\mathcal{J}_{20},\qquad
\mathcal{K}_3=\mathcal{J}_{30}.
\tag{5.6.3}\label{eq:5.6.3}
\end{equation}
Eq.~\eqref{eq:5.6.1} then reads
\begin{equation}
[\mathcal{J}_i,\mathcal{J}_j]=i\epsilon_{ijk}\mathcal{J}_k,
\tag{5.6.4}\label{eq:5.6.4}
\end{equation}
\begin{equation}
[\mathcal{J}_i,\mathcal{K}_j]=i\epsilon_{ijk}\mathcal{K}_k,
\tag{5.6.5}\label{eq:5.6.5}
\end{equation}
\begin{equation}
[\mathcal{K}_i,\mathcal{K}_j]=-i\epsilon_{ijk}\mathcal{J}_k,
\tag{5.6.6}\label{eq:5.6.6}
\end{equation}
where $i,j,k$ run over the values $1,2,3$, and $\epsilon_{ijk}$ is the totally
antisymmetric quantity with $\epsilon_{123}=+1$. Eq.~\eqref{eq:5.6.4} just says that
the matrices $\mathcal{J}$ generate a representation of the rotation subgroup
of the Lorentz group, and Eq.~\eqref{eq:5.6.5} just represents the fact that
$\mathcal{K}$ is a three-vector. The minus sign in the right-hand side of Eq.
\eqref{eq:5.6.6} arises from the fact that $\eta_{00}=-1$, and plays a crucial role in
what follows.

It is very convenient to replace the matrices $\mathcal{J}$ and $\mathcal{K}$
with two decoupled spin three-vectors, writing
\begin{equation}
\mathcal{A}\equiv\frac12(\mathcal{J}+i\mathcal{K}),
\tag{5.6.7}\label{eq:5.6.7}
\end{equation}
\begin{equation}
\mathcal{B}\equiv\frac12(\mathcal{J}-i\mathcal{K}).
\tag{5.6.8}\label{eq:5.6.8}
\end{equation}
It is easy to see that the commutation relations \eqref{eq:5.6.4}--\eqref{eq:5.6.6} are
equivalent to
\begin{equation}
[\mathcal{A}_i,\mathcal{A}_j]=i\epsilon_{ijk}\mathcal{A}_k,
\tag{5.6.9}\label{eq:5.6.9}
\end{equation}
\begin{equation}
[\mathcal{B}_i,\mathcal{B}_j]=i\epsilon_{ijk}\mathcal{B}_k,
\tag{5.6.10}\label{eq:5.6.10}
\end{equation}
\begin{equation}
[\mathcal{A}_i,\mathcal{B}_j]=0.
\tag{5.6.11}\label{eq:5.6.11}
\end{equation}
We find matrices satisfying Eqs.~\eqref{eq:5.6.9}--\eqref{eq:5.6.11} in the same way that we
find matrices representing the spins of a pair of uncoupled particles---as a
direct sum. That is, we label the rows and columns of these matrices with a
pair of integers and/or half-integers $a,b$, running over the values
\begin{equation}
a=-A,-A+1,\ldots,+A,
\tag{5.6.12}\label{eq:5.6.12}
\end{equation}
\begin{equation}
b=-B,-B+1,\ldots,+B
\tag{5.6.13}\label{eq:5.6.13}
\end{equation}
and take\footnote{There is an alternative formalism~\hyperref[ref:5.8]{[8]}, based on the fact that
the spin $j$ representation of the rotation group can be written as the
symmetrized direct product of $2j$ spin $1/2$ representations---i.e., as a
symmetric $\mathrm{SU}(2)$ tensor with $2j$ two-valued indices. We can
therefore write fields belonging to the $(A,B)$ representation with $2A$
two-valued $(1/2,0)$ indices and $2B$ two-valued $(0,1/2)$ indices, the latter
written with dots to distinguish them from the former.}
\begin{equation}
(\mathcal{A})_{a'b',ab}=\delta_{b'b}\mathbf J^{(A)}_{a'a},
\tag{5.6.14}\label{eq:5.6.14}
\end{equation}
\begin{equation}
(\mathcal{B})_{a'b',ab}=\delta_{a'a}\mathbf J^{(B)}_{b'b},
\tag{5.6.15}\label{eq:5.6.15}
\end{equation}
where $\mathbf J^{(A)}$ and $\mathbf J^{(B)}$ are the standard spin matrices
for spins $A$ or $B$:
\begin{equation}
(\mathbf J_3^{(A)})_{a'a}=a\delta_{a'a},
\tag{5.6.16}\label{eq:5.6.16}
\end{equation}
\begin{equation}
(\mathbf J_1^{(A)}\mathbin{\pm}i\mathbf J_2^{(A)})_{a'a}
=\delta_{a',a\pm1}\sqrt{(A\mp a)(A\pm a+1)},
\tag{5.6.17}\label{eq:5.6.17}
\end{equation}
and likewise for $\mathbf J^{(B)}$. The representation is labelled by the
values of the positive integers and/or half-integers $A$ and $B$. We see that
the $(A,B)$ representation has dimensionality $(2A+1)(2B+1)$.

The finite-dimensional representations of the homogeneous Lorentz group are
not unitary, because $\mathcal{A}$ and $\mathcal{B}$ are Hermitian, and
therefore $\mathcal{J}$ is Hermitian but $\mathcal{K}$ is
\emph{anti-Hermitian}. This is because of the $i$ in Eqs.~\eqref{eq:5.6.7} and \eqref{eq:5.6.8},
which is required by the minus sign in \eqref{eq:5.6.6}, and hence stems from the fact
that the homogeneous Lorentz group is not the same as the four-dimensional
rotation group $\mathrm{SO}(4)$, a compact group, but instead is the
non-compact group known as $\mathrm{SO}(3,1)$. It is only compact groups that
can have finite-dimensional unitary representations (aside from
representations in which the non-compact part is represented trivially, by the
identity). There is no problem in working with non-unitary representations,
because the objects we are now concerned with are fields, not wave functions,
and do not need to have a Lorentz-invariant positive norm.

In contrast, the rotation group is represented unitarily, with its generators
represented by the Hermitian matrices
\begin{equation}
\mathcal{J}=\mathcal{A}+\mathcal{B}.
\tag{5.6.18}\label{eq:5.6.18}
\end{equation}
By the usual rules of vector addition, we can see that a field that transforms
according to the $(A,B)$ representation of the homogeneous Lorentz group has
components that rotate like objects of spin $j$, with
\[
j=A+B,A+B-1,\ldots,|A-B|.
\]
This is enough to identify the $(A,B)$ representations with the perhaps more
familiar tensors and spinors. For instance, a $(0,0)$ field is obviously
scalar, with only a single $j=0$ component. A $(1/2,0)$ or $(0,1/2)$ field can
only have $j=1/2$; these are the top (i.e., $\gamma_5=+1$) and bottom
($\gamma_5=-1$) two components of the Dirac spinor. A $(1/2,1/2)$ field has
components with $j=1$ and $j=0$, corresponding to the spatial part
$\mathbf v$ and time-component $v^0$ of a four-vector $v^\mu$. More generally,
an $(A,A)$ field contains terms with only integer spins
$2A,2A-1,\ldots,0$, and corresponds to a traceless symmetric tensor of rank
$2A$. (Note that the number of independent components of a symmetric tensor
of rank $2A$ in four dimensions is
\[
\frac{4\cdot5\cdots(4+2A-1)}{(2A)!}=\frac{(3+2A)!}{6(2A)!}
\]
and the tracelessness condition reduces this to
\[
\frac{(3+2A)!}{6(2A)!}-\frac{(1+2A)!}{6(2A-2)!}=(2A+1)^2
\]
as expected for an $(A,A)$ field.) One more example: a $(1,0)$ or $(0,1)$
field can only have $j=1$, and corresponds to an antisymmetric tensor
$F^{\mu\nu}$ that satisfies the further irreducibility `duality' conditions
\[
F^{\mu\nu}=\mathbin{\pm}\frac{i}{2}\epsilon^{\mu\nu\lambda\rho}F_{\lambda\rho}
\]
for $(1,0)$ and $(0,1)$ fields, respectively. Of course, it is only in four
dimensions that an antisymmetric two-index tensor $F^{\mu\nu}$ can be divided
into such `self-dual' and `anti-self-dual' parts.

A general tensor of rank $N$ transforms as the direct product of $N$
$(1/2,1/2)$ four-vector representations. It can therefore be decomposed (by
suitable symmetrizations and antisymmetrizations and extracting traces) into
irreducible terms $(A,B)$ with $A=N/2,N/2-1,\ldots$ and
$B=N/2,N/2-1,\ldots$. In this way, we can construct any irreducible
representation $(A,B)$ for which $A+B$ is an integer. The spin
representations, for which $A+B$ is half an odd integer, can similarly be
constructed from the direct product of these tensor representations and the
Dirac representation $(1/2,0)\oplus(0,1/2)$.

For instance, taking the direct product of the vector $(1/2,1/2)$
representation and the Dirac $(1/2,0)\oplus(0,1/2)$ representation gives a
spinor-vector $\psi^\mu$, that transforms according to the reducible
representation
\[
(\tfrac12,\tfrac12)\otimes[(\tfrac12,0)\oplus(0,\tfrac12)]
=(\tfrac12,1)\oplus(\tfrac12,0)\oplus(1,\tfrac12)\oplus(0,\tfrac12).
\]
The quantity $\gamma_\mu\psi^\mu$ would transform as an ordinary
$(1/2,0)\oplus(0,1/2)$ Dirac field, so we can isolate the
$(1/2,1)\oplus(1,1/2)$ representation\footnote{According to Eq.~\eqref{eq:5.6.18},
such a field transforms under ordinary rotations as a direct sum of two
$j=3/2$ and two $j=1/2$ components. The doubling is eliminated by imposing
the Dirac equation $[\gamma^\nu\partial_\nu+m]\psi^\mu=0$, and the remaining
$j=1/2$ component is eliminated by requiring that $\partial_\mu\psi^\mu=0$.
With these conditions, the field describes a single particle of spin $j=3/2$.}
by requiring that $\gamma_\mu\psi^\mu=0$. This is the
\emph{Rarita--Schwinger field}~\hyperref[ref:5.9]{[9]}.

So far in this section we have only considered the representations of the
proper orthochronous Lorentz group. In any representation of the Lorentz group
including space inversion, there must be a matrix $\beta$ which reverses the
signs of tensors with odd numbers of space indices, and in particular
\begin{equation}
\beta\mathcal{J}\beta^{-1}=+\mathcal{J},
\qquad
\beta\mathcal{K}\beta^{-1}=-\mathcal{K}.
\tag{5.6.19}\label{eq:5.6.19}
\end{equation}
In terms of the matrices \eqref{eq:5.6.7} and \eqref{eq:5.6.8}, this is
\begin{equation}
\beta\mathcal{A}\beta^{-1}=\mathcal{B},
\qquad
\beta\mathcal{B}\beta^{-1}=\mathcal{A}.
\tag{5.6.20}\label{eq:5.6.20}
\end{equation}
Thus an irreducible $(A,B)$ representation of the proper orthochronous
homogeneous Lorentz group does not provide a representation of the Lorentz
group including space inversion unless $A=B$. As we have seen, these $(A,A)$
representations are the scalar, the vector, and the symmetric traceless
tensors. For $A\ne B$, the irreducible representations of the Lorentz group
including space inversion are the direct sums $(A,B)\oplus(B,A)$, of
dimensionality $2(2A+1)(2B+1)$. One of these is the
$(1/2,0)\oplus(0,1/2)$ Dirac representation discussed in Section 5.4. The
$4\times4$ matrix \eqref{eq:5.4.29} provides the $\beta$-matrix for this representation.
Another familiar example is the $(1,0)\oplus(0,1)$ representation, which as we
have seen is just the antisymmetric tensor of second rank, including both
self-dual and anti-self-dual parts.
